We compare our graph-guided prompting workflow against a set of baseline and ablation pipelines that reflect a standard, linear section-generation approach. The principal baseline—henceforth the ``linear prompting'' pipeline—generates document sections in a sequential, heading-by-heading manner without an explicit document graph or scheduler. Each section is produced by a single prompting call that supplies the model with the draft document context (previously generated sections and the global outline), any retrieved evidence made available to that call, and instructions to emit section text. This baseline is intended to represent a straightforward application of retrieval-augmented generation and retrieval-aware sectioning; prior work has documented a variety of behaviors and failure modes in retrieval pipelines that motivate careful, controlled baseline construction \cite{web_penha_2021_evaluating_the_robustness_of_retrieval_pipelines,web_penha_2023_do_the_findings_of_document_and_passage_retrieva}. 

To isolate the contributions of individual design choices, we evaluate a small set of controlled baseline variants and ablations:

\begin{itemize}
  \item Prompt-order variants: (i) top-down linear ordering, where sections are generated following the outline from the highest-level headings downward; and (ii) leaf-first ordering, where leaf sections are generated earlier and higher-level sections are composed from the aggregated leaf outputs. These variants probe the extent to which generation order alone affects citation alignment and coherence.
  \item Evidence-availability ablations: (i) retrieval-enabled, in which each section prompt is supplied with the same retrieval corpus and retrieved passages; and (ii) retrieval-disabled, where prompts do not receive external evidence and must rely solely on model-internal knowledge. This contrast measures the impact of explicit evidence injection independent of scheduling and is informed by known sensitivities of retrieval pipelines to query and corpus variation \cite{web_penha_2021_evaluating_the_robustness_of_retrieval_pipelines}.
  \item Output-constraint ablations: (i) JSON-constrained generation, enforcing the same strict section-level schema used in the graph-guided pipeline; and (ii) free-text generation, in which the model is only asked for plain LaTeX or paragraph text. This pair isolates the effect of structured, machine-parseable outputs on citation placement and downstream assembly.
  \item Scheduler ablation: a hybrid pipeline that retains explicit retrieval and JSON constraints but omits the dependency-respecting scheduler, producing sections in a randomized or outline-ordered sequence. This tests whether scheduling (the primary novelty of the graph-guided approach) is necessary for observed gains.
\end{itemize}

Across all baselines and ablations we hold constant implementation factors that are not under test: the underlying language model interface, the retrieval corpus and index, and the core prompt templates insofar as they are comparable (for example, retrieval prompts include the same retrieved passage summaries when retrieval is enabled). Maintaining these controls ensures that differences attributed to the graph-guided pipeline derive from the document-graph representation, the scheduling algorithm, and the section-level evidence-injection mechanism rather than from model capacity or retrieval coverage.

General background knowledge: linear prompting and retrieval-augmented generation are common starting points for multi‑section generation experiments; the above baseline variants therefore reflect standard axes of methodological variation (prompt order, access to retrieved evidence, and output structure). Finally, the chosen baselines and ablations are intended to be conservative and interpretable—each is designed to isolate a single factor (order, evidence, constraints, scheduler) so that any improvements observed with the graph-guided workflow can be more confidently attributed to its graph-aware components. Limitations and suggestions for additional baselines (e.g., multi‑pass planners or hierarchical latent-variable generators) are discussed in the Limitations and Future Work subsections.